\chapter{چکیده}
  
در عصر تحولات شتابان فناوریهای خودران، سیستمهای بینایی ماشین به عنوان یکی از ارکان حیاتی در ادراک محیطی خودروها نقش ایفا میکنند. پروژه حاضر با تمرکز بر بهبود عملکرد سیستمهای تشخیص علائم راهنمایی و رانندگی، تلاشی هدفمند برای تقریب هرچه بیشتر رفتار این سیستمها به قابلیتهای انسانی است. اهمیت این موضوع در شرایط واقعی رانندگی خودمختار، که کوچکترین خطا در تشخیص علائم میتواند پیامدهای جبرانناپذیری داشته باشد، بیش از پیش آشکار میشود. در این راستا، پژوهش ما با بهرهگیری از مدلهای پیشرفته هوش مصنوعی و بینایی کامپیوتر، رویکردی تحلیلی-معنایی را برای درک عمیقتر ویژگیهای علائم راهنمایی و رانندگی طراحی کرده است. 
مطالعه گسترده و نظاممندِ پژوهشهای پیشین، نخستین گام این پروژه بود. با بررسی انتقادی کارهای مرتبط، نقاط قوت و ضعف روشهای موجود شناسایی شد و این تحلیل به انتخاب هوشمندانهترین مجموعه داده متناسب با اهداف و محدودیتهای پروژه منجر گردید. در گام بعدی، با تمرکز بر درک جامعِ مجموعه داده انتخابی، ویژگیهای کلیدی علائم از جمله شکل، رنگ، نمادها، و زمینه قرارگیری آنها در محیط واقعی مورد تحلیل قرار گرفتند. رویکرد این پژوهش صرفاً به حجم انبوه دادهها متکی نبود، بلکه با تأکید بر تحلیل معنایی و ارتباط بین ویژگیهای ساختاری علائم (مانند تقارن، تضاد رنگی، و موقعیت مکانی) و نقش آنها در انتقال پیام به رانندگان، لایهای جدید از هوشمندی را به سیستم افزود. 
یکی از نوآوریهای اصلی این پروژه، شبیهسازی شرایط پیچیده و غیرایدهآل محیطی است که علائم در آن قرار میگیرند. با استفاده از ابزارهای تولید دادههای مصنوعی (\LTRfootnote{Synthetic Data}) و تکنیکهای افزایش داده (\LTRfootnote{Data Augmentation}), سناریوهایی مانند تغییرات نورپردازی (\LTRfootnote{Night/Day, Sun Reflections}), اعوجاجهای دیداری (\LTRfootnote{Rain, Snow, Fog}), و حتی آسیبهای فیزیکی به علائم (\LTRfootnote{Scratches, Dirt}) مدلسازی شدند. این شبیهسازیها نه تنها به بهبود مقاومت (\LTRfootnote{Robustness}) مدل در مواجهه با دادههای واقعی کمک کرد، بلکه امکان ارزیابی عملکرد سیستم در حالات مرزی (\LTRfootnote{Edge Cases}) را نیز فراهم آورد. 
در نهایت، مدل پیشنهادی با تلفیق رویکردهای مبتنی بر یادگیری عمیق (\LTRfootnote{Deep Learning}) (مانند شبکه های کانولوشنی) و پردازشهای معنایی لایه لایه، به سمت سیستمیک سوق داده شد که نه تنها قادر به تشخیص دقیق علائم است، بلکه میتواند با استنتاج منطقی از جایگاه علائم در چارچوب قوانین راهنمایی و رانندگی (\LTRfootnote{Traffic Rules}) (مانند اولویتبندی علائم متضاد), گامی به سوی "استدلال انسانی" بردارد. ارزیابی نتایج نشان میدهد که این سیستم در مقایسه با روشهای متکی صرف بر حجم داده، از دقت تشخیص بالاتری در شرایط پویا و غیرقابل پیشبینی برخوردار است. این پژوهش مسیری نوین در همگرایی هوش مصنوعی (\LTRfootnote{AI}), مهندسی معنایی, و الزامات امنیتی سیستمهای خودران (\LTRfootnote{Autonomous Systems}) ترسیم میکند و چارچوبی را برای توسعه سیستمهای تشخیص علائم با قابلیت تفسیرپذیری (\LTRfootnote{Interpretability}) بالا ارائه مینماید. 
  
  
  
  